\subsection{Staging und Production Umgebung}

Die Docker-Compose-Dateien sind im Unterordner /.deployment zu finden.
In der Produktionsumgebung läuft die aktuelle Live-Version von Codetask, während die Staging-Umgebung einen Zwischenstand zwischen der lokalen und der Produktionsumgebung darstellt.
In dieser können Implementierungen und Entwicklungen in einer live-ähnlichen Umgebung getestet werden.

Die Compose-Dateien ähneln denen der lokalen Umgebung, haben dennoch deutliche Unterschiede.
So gibt es beispielsweise einen weiteren Service namens Traefik, der als Reverse Proxy benutzt wird. 
Die Umgebungen haben auch einen Eintrag des DNS-Servers der HTWG, wodurch sie auch von außen erreichbar sind. Die Staging-Umgebung ist hingegen nur über VPN erreichbar.
Der Host der Prod-Umgebung ist codetask.in.htwg-konstanz.de, der Host der Staging-Umgebung ist staging-codetask.in.htwg-konstanz.de. Wird der Host für die API oder AI-API verwendet, wird jeweils das Passende vorne dran gehängt, z. B. „api.codetask.in.htwg-konstanz.de”.
Daher haben beide Umgebungen jeweils drei URLs, wodurch drei Services von außen erreichbar sind. Die restlichen Services sind nur intern im Compose-Netzwerk erreichbar.
Bei der Produktionsumgebung haben alle drei URLs ein TLS-Zertifikat erhalten, wodurch jede mit HTTPS genutzt werden kann. Bei der Staging-Umgebung hat nur die URL für das Webframework ein TLS-Zertifikat erhalten.

Anders als bei der lokalen Umgebung werden die Images bei beiden Compose nicht mehr mit Dockerfiles selbst gebaut, sondern über die Container Registry vom GitLab-Repository von Codetask geholt.
Die Staging- und Produktionsumgebung haben auch ein anderes Dockerfile als die lokale Umgebung. Ein Unterschied ist beispielsweise, dass die Container nicht im DEV-Modus laufen oder dass im Frontend-Dockerfile ein Container für das Webframework und den Webserver enthalten ist.